% ===========================================================================
% Category: Self-Organizing Map
% Generated from ML_Lab_Manual_Completed.md - edit the Markdown and re-run
% scripts/build_latex_manual.py instead of editing this file by hand.
% ===========================================================================

\chapter{Experiment 15: Self-Organizing Map
(SOM)}\label{experiment-15-self-organizing-map-som}

\textbf{Category:} Unsupervised neural learning

\section{1. Objective}\label{objective}

Map high-dimensional observations onto a low-dimensional grid while
preserving neighborhood structure; evaluate with quantization and
topographic error.

\section{2. Dataset Used}\label{dataset-used}

UCI Wine: 178 samples, 13 chemical features, 3 cultivars (labels used
only for visualization). MinMax-scaled to {[}0,1{]}.

\section{3. Required Libraries}\label{required-libraries}

minisom, numpy, pandas, matplotlib

\section{4. Brief Theory}\label{brief-theory}

A SOM has a grid of neurons with weight vectors. For each input the Best
Matching Unit (BMU) is found (minimum distance), then the BMU and its
Gaussian neighborhood are moved toward the input: w\_j ← w\_j + α(t)
h\_cj(t) (x − w\_j), with decaying α(t) and σ(t). The result is a
topology-preserving 2D projection; the U-Matrix visualizes cluster
boundaries.

\section{5. Procedure}\label{procedure}

\begin{enumerate}
\def\labelenumi{\arabic{enumi}.}
\tightlist
\item
  Load wine; MinMax scale.
\item
  Create a 12×12 SOM with Gaussian neighborhood (σ=1.5, lr=0.5);
  initialize weights with PCA.
\item
  Train for 5,000 random samples.
\item
  Compute quantization error (average sample→BMU distance) and
  topographic error (fraction of non-adjacent top-2 BMUs).
\item
  Plot the U-Matrix with each sample projected onto its BMU (marker =
  true cultivar).
\item
  Discuss grid size and σ.
\end{enumerate}

\section{6. Reference Implementation}\label{reference-implementation}

\begin{lstlisting}[language=Python]
som = MiniSom(x=12, y=12, input_len=13, sigma=1.5, learning_rate=0.5,
              neighborhood_function='gaussian', random_seed=42)
som.pca_weights_init(X)
som.train_random(data=X, num_iteration=5000)
qe = som.quantization_error(X); te = som.topographic_error(X)
\end{lstlisting}

\section{7. Results}\label{results}

{\def\LTcaptype{none} % do not increment counter
\begin{longtable}[]{@{}ll@{}}
\toprule\noalign{}
Metric & Value \\
\midrule\noalign{}
\endhead
\bottomrule\noalign{}
\endlastfoot
Initial quantization error & 0.4447 \\
Final quantization error & \textbf{0.1965} \\
Topographic error & \textbf{0.0393} (\textasciitilde4\% foldings) \\
\end{longtable}
}

On the U-Matrix, the three wine cultivars occupy contiguous regions
separated by bright ridges, even though the SOM never saw the labels ---
evidence of unsupervised structure discovery.

\section{8. Discussion and Conclusion}\label{discussion-and-conclusion}

The map halved its quantization error during training and preserves
topology well (topographic error \textless{} 0.05). PCA initialization
gave a well-unfolded map from the start. A larger grid reduces
quantization error but risks fragmenting clusters; σ/lr must decay for
stable convergence.

\section{9. Answers to Viva Questions}\label{answers-to-viva-questions}

\begin{itemize}
\tightlist
\item
  \textbf{SOM vs K-means?} Both compress data by prototypes, but K-means
  only finds centroids while SOM arranges them on a grid so that
  neighboring neurons represent similar inputs (topology preservation).
\item
  \textbf{What is a BMU?} The neuron whose weight vector is closest to
  the input, i.e.~the winner of the competition.
\item
  \textbf{Why is SOM called a topology-preserving map?} Because the
  update neighborhood ensures nearby grid positions end up with similar
  weight vectors, so high-dimensional neighborhoods are approximately
  preserved in 2D.
\end{itemize}

\begin{center}\rule{0.5\linewidth}{0.5pt}\end{center}

\section{10. Complete Code Listing}\label{complete-code-listing}

\emph{Full runnable program for this experiment, extracted from
\passthrough{\lstinline!notebooks/07\_self\_organizing\_map.ipynb!}
(executed; results above are its actual output).}

\textbf{Listing 15.1 --- Core numerical and plotting libraries}

\begin{lstlisting}[language=Python]
# ---- Core numerical and plotting libraries ----
import numpy as np
import pandas as pd
import matplotlib.pyplot as plt
import seaborn as sns

# ---- Dataset + scaling ----
from sklearn.datasets import load_wine
from sklearn.preprocessing import MinMaxScaler
# ---- MiniSom: lightweight Kohonen Self-Organizing Map implementation ----
from minisom import MiniSom

# Styling setup (consistent look across all notebooks)
plt.style.use('seaborn-v0_8-whitegrid' if 'seaborn-v0_8-whitegrid' in plt.style.available else 'default')
plt.rcParams['font.sans-serif'] = 'DejaVu Sans'
plt.rcParams['figure.figsize'] = (10, 6)
plt.rcParams['figure.dpi'] = 110

np.random.seed(42)
print("MiniSom and libraries loaded successfully!")

# ---- Load the Wine recognition dataset ----
wine = load_wine()
X_raw = wine.data          # 178 samples x 13 chemical features
y = wine.target            # 3 wine cultivars (only used for visualization - SOM is unsupervised)
feature_names = wine.feature_names
target_names = wine.target_names

print(f"Samples: {X_raw.shape[0]}, Features: {X_raw.shape[1]}, Classes: {target_names.tolist()}")

# ---- MinMax scaling to [0, 1] ----
# SOM weight vectors live in the same range as the inputs, so scaling keeps learning stable.
scaler = MinMaxScaler()
X = scaler.fit_transform(X_raw)
\end{lstlisting}

\textbf{Listing 15.2 --- Configure the Kohonen grid}

\begin{lstlisting}[language=Python]
# ---- Configure the Kohonen grid ----
grid_x, grid_y = 12, 12  # 12x12 grid = 144 neurons (fewer than the 178 samples)
input_len = X.shape[1]   # 13 chemical features per neuron weight vector

som = MiniSom(
    x=grid_x,
    y=grid_y,
    input_len=input_len,
    sigma=1.5,                       # initial neighborhood radius
    learning_rate=0.5,               # initial learning rate
    neighborhood_function='gaussian',# smooth Gaussian neighborhood kernel
    random_seed=42
)

# ---- Initialize weights with PCA for topologically faithful, faster convergence ----
som.pca_weights_init(X)
print("Initial Quantization Error:", som.quantization_error(X))

# ---- Train: 5000 random samples presented to the map ----
som.train_random(data=X, num_iteration=5000, verbose=False)
final_qe = som.quantization_error(X)   # average distance from samples to their BMU
final_te = som.topographic_error(X)    # fraction of samples whose 2nd BMU is not a grid neighbor

print(f"Training complete after 5,000 iterations!")
print(f" - Final Quantization Error (mean distance to BMU): {final_qe:.4f}")
print(f" - Final Topographic Error (proportion of foldings): {final_te:.4f}")
\end{lstlisting}

\textbf{Listing 15.3 --- Plot the U-Matrix (unified distance map)}

\begin{lstlisting}[language=Python]
# ---- Plot the U-Matrix (unified distance map) ----
# Each cell = average distance between a neuron and its immediate grid neighbors.
# Dark valleys = dense clusters; bright ridges = boundaries between clusters.
plt.figure(figsize=(10, 8))
plt.pcolor(som.distance_map().T, cmap='bone_r', alpha=0.9)
cbar = plt.colorbar()
cbar.set_label("Normalized Inter-Neuron Distance (U-Matrix)", fontsize=11)

# ---- Project every sample onto its Best Matching Unit (BMU) ----
markers = ['o', 's', '^']
colors = ['red', 'green', 'blue']

for idx, x_vec in enumerate(X):
    w = som.winner(x_vec)  # grid coordinates of the best matching neuron for this sample
    plt.plot(
        w[0] + 0.5,        # +0.5 centers the marker inside the pcolor cell
        w[1] + 0.5,
        markers[y[idx]],   # marker shape encodes the true cultivar
        markerfacecolor='None',
        markeredgecolor=colors[y[idx]],
        markersize=10,
        markeredgewidth=2
    )

# Build a legend for the three cultivars
for c_idx, c_name in enumerate(target_names):
    plt.plot([], [], marker=markers[c_idx], color=colors[c_idx], linestyle='None',
             label=f'Cultivar {c_name}', markersize=9, markeredgewidth=2)

plt.title("SOM U-Matrix with Projected Wine Cultivars", fontsize=14, fontweight='bold')
plt.legend(bbox_to_anchor=(1.25, 1), loc='upper left')
plt.tight_layout()
plt.show()
\end{lstlisting}
